\section{Основні означення}

У цьому розділі наведено основні поняття, означення та теоретичні
положення~\cite{jurafsky}, які використовуються в подальшому викладі
кваліфікаційної роботи. Розглядаються
базові складові систем автоматичного перекладу мовлення, формальні постановки
відповідних задач та метрики оцінювання якості перекладу.

\subsection{Мовлення як сигнал та його подання}

\textit{Мовлення} --- це нестаціонарний акустичний сигнал, що формується в результаті
артикуляційної діяльності людини та несе лінгвістичну інформацію. У цифрових
системах мовлення подається у вигляді дискретного сигналу
\[
x[n] = x(nT_s),
\]
де $T_s$ --- період дискретизації, що визначається частотою дискретизації
$f_s = 1/T_s$.

Для обробки мовлення застосовують \textit{спектральні ознаки}, отримані за допомогою
короткочасного перетворення Фур'є (STFT). Одними з найпоширеніших ознак є
мел-частотні кепстральні коефіцієнти (MFCC) та мел-спектрограми, які
узгоджуються з особливостями слухового сприйняття людини.

\subsection{Автоматичне розпізнавання мовлення}
\label{subsec:asr}

\textit{Автоматичне розпізнавання мовлення} (Automatic Speech Recognition, ASR) ---
це задача перетворення мовного аудіосигналу у текстову послідовність. Формально
ASR подають як задачу знаходження найбільш імовірної текстової послідовності
$y$ за наявності акустичного сигналу $x$:
\[
\hat{y} = \arg\max_{y} P(y \mid x).
\]

Сучасні ASR-системи базуються на нейронних мережах глибинного
навчання~\cite{asr_survey,goodfellow_dl}, зокрема на згорткових,
рекурентних та трансформерних архітектурах. Широко
застосовуються end-to-end підходи, які безпосередньо моделюють умовний розподіл
$P(y \mid x)$ без явного поділу на акустичну та мовну моделі.

Для мов із багатою морфологією, зокрема української, важливими є використання
субслівних одиниць (byte-pair encoding, unigram language models), що дозволяє
зменшити словникову розмірність.

Основною метрикою якості ASR є \textit{Word Error Rate} (WER):
\[
\text{WER} = \frac{S + D + I}{N},
\]
де $S$ --- кількість замін, $D$ --- кількість вилучень, $I$ --- кількість вставок,
а $N$ --- загальна кількість слів у референсному тексті. Для морфологічно багатих
мов додатково використовують \textit{Character Error Rate} (CER), що обчислюється
аналогічно на рівні символів.

\subsection{Машинний переклад}

\textit{Машинний переклад} (Machine Translation, MT) --- це задача автоматичного
перетворення тексту з мови джерела на мову призначення зі збереженням
семантичного змісту. У нейронному машинному перекладі задача формалізується як
\[
\hat{z} = \arg\max_{z} P(z \mid y),
\]
де $y$ --- вхідна текстова послідовність мовою джерела, а $z$ --- відповідна
послідовність мовою призначення.

Найбільш поширеною архітектурою для MT є \textit{Transformer} \cite{vaswani}, який
ґрунтується на механізмі самоуваги (self-attention). Перевагою таких моделей є
здатність моделювати довгі залежності у тексті та виконувати паралельні
обчислення.

У контексті автоматичного перекладу мовлення в реальному часі значення мають
моделі одночасного перекладу (simultaneous translation), які генерують переклад
частинами ще до отримання повного вхідного речення.

\subsection{Синтез мовлення}

\textit{Синтез мовлення} (Text-to-Speech, TTS) --- це процес перетворення текстової
інформації у штучно згенероване мовлення. TTS є третім компонентом повного
конвеєра перекладу мовлення (ASR $\to$ MT $\to$ TTS), проте в даному
дослідженні не використовується, оскільки робота зосереджена на оцінюванні
якості текстового перекладу.

\subsection{Каскадна архітектура перекладу мовлення}
\label{subsec:cascade}

\textit{Автоматичний переклад мовлення} (Speech Translation, ST) --- це процес
перетворення мовлення однією мовою на текст або мовлення іншою мовою.
Системи перекладу мовлення можуть працювати в офлайн-режимі (обробка
попередньо записаного аудіо) або в режимі, наближеному до реального
часу~\cite{real_time}. У даному дослідженні розглядається офлайн-режим,
що дозволяє зосередитися на аналізі якості перекладу без обмежень
латентності.

Каскадну архітектуру перекладу мовлення можна подати у вигляді
послідовного конвеєра:
\[
x(t) \xrightarrow{\ \text{ASR}\ } \hat{y} \xrightarrow{\ \text{MT}\ } \hat{z},
\]
де $x(t)$ --- вхідний мовний сигнал, $\hat{y}$ --- розпізнана текстова
послідовність мовою джерела, $\hat{z}$ --- перекладений текст мовою
призначення. Для повного конвеєра може додаватися етап синтезу мовлення
(TTS): $\hat{z} \xrightarrow{\text{TTS}} \hat{s}(t)$, проте в даному
дослідженні оцінюється лише текстовий переклад (етапи ASR + MT).

Ключовою особливістю каскадної архітектури є \textit{каскадний ефект
поширення помилок}~\cite{error_propagation}: помилки, допущені на
етапі ASR, потрапляють на вхід MT-компоненту та можуть суттєво
погіршити якість фінального перекладу.

\subsection{Семантична точність перекладу}

\textit{Семантична точність} --- це ступінь збереження змісту оригінального
висловлювання у перекладі. Оцінювання семантичної точності здійснюється за
допомогою автоматичних метрик та шляхом експертного аналізу.

Якість перекладу оцінюється на текстовому рівні --- порівнюється результат
перекладу $\hat{z}$ із референсним перекладом $z^*$. Автоматичні метрики
формалізуються як функції подібності між цими послідовностями:
\[
Q = f(\hat{z},\, z^*),
\]
де $Q$ --- числовий показник якості, а $f$ --- відповідна метрика.

\textit{BLEU} (Bilingual Evaluation Understudy) \cite{papineni_bleu} --- метрика,
що базується на порівнянні $n$-грам перекладу з референсним текстом та є
стандартом де-факто для оцінювання машинного перекладу.

\textit{BERTScore} \cite{bert_score} --- метрика, що використовує контекстуальні
векторні подання мовної моделі BERT для обчислення семантичної подібності.
На відміну від BLEU, BERTScore здатна виявляти семантичну еквівалентність
за відсутності лексичного збігу.

\textit{METEOR} \cite{banerjee_meteor} --- метрика, що враховує не лише точний
збіг слів, а й стемінг, синоніми та порядок слів.

Детальний опис формул та обґрунтування вибору метрик наведено в
підрозділі~\ref{subsec:metrics}.

\subsection{Міжфразова когерентність}

\textit{Міжфразова когерентність} --- це ступінь контекстуальної зв'язності між
послідовними сегментами перекладу. У каскадних системах перекладу мовлення кожен
сегмент обробляється незалежно, що означає, що система працює з
обмеженим контекстом, що може призводити до порушення зв'язності: вибір
значення полісемічного слова у попередньому сегменті може не узгоджуватися з
контекстом наступного, займенникові посилання можуть втрачати антецедент, а
тематичні переходи --- залишатися непозначеними.

Когерентність між $i$-м та $(i+1)$-м сегментами перекладу можна оцінити як
\[
C_i = g(\hat{z}_i,\, \hat{z}_{i+1},\, z^*_i,\, z^*_{i+1}),
\]
де $g$ --- функція, що враховує збереження референтних зв'язків, тематичної
послідовності та контекстуальної узгодженості між сусідніми сегментами.

Порушення когерентності є критичним для слухача, оскільки навіть за умови
високої семантичної точності окремих сегментів втрата зв'язності між ними може
суттєво ускладнити розуміння загального змісту повідомлення.

\subsection{Класифікація помилок перекладу}
\label{subsec:error_class}

Для детального аналізу якості перекладу застосовують класифікацію помилок за
типами:
\begin{itemize}[nosep]
    \item \textit{Пропуски} (omissions) --- частина змісту оригіналу відсутня у
    перекладі;
    \item \textit{Спотворення} (distortions) --- зміст оригіналу передано невірно;
    \item \textit{Додавання} (additions) --- у перекладі присутня інформація,
    якої немає в оригіналі.
\end{itemize}

Для кількісного оцінювання обчислюють відповідні коефіцієнти:
\[
R_{\text{omit}} = \frac{N_{\text{omit}}}{N_{\text{total}}}, \quad
R_{\text{dist}} = \frac{N_{\text{dist}}}{N_{\text{total}}}, \quad
R_{\text{add}} = \frac{N_{\text{add}}}{N_{\text{total}}},
\]
де $N_{\text{omit}}, N_{\text{dist}}, N_{\text{add}}$ --- кількість сегментів із
відповідним типом помилки, а $N_{\text{total}}$ --- загальна кількість сегментів.

\subsection{Огляд літератури}

Дослідження автоматичного перекладу мовлення є активним напрямком
на перетині обробки мовлення та машинного перекладу. Сперберта та
Паулік~\cite{cascade_st} надають ґрунтовний порівняльний аналіз
каскадних та end-to-end підходів до перекладу мовлення, зазначаючи, що
каскадні системи залишаються конкурентоспроможними завдяки модульності
та можливості покомпонентного оцінювання. Анастасопулос та
Чань~\cite{st_survey} у своєму огляді систематизують досягнення
end-to-end моделей, але підкреслюють, що для більшості мовних пар
каскадні системи забезпечують вищу якість.

\paragraph{Каскадний ефект поширення помилок.}
Центральним питанням для каскадних систем є вплив помилок ASR на якість
фінального перекладу. Руіз та Федеріко~\cite{error_propagation}
експериментально показали, що підвищення WER на 10\% призводить до
зниження BLEU на 5--15 пунктів залежно від мовної пари та домену.
Бентіволі та ін.~\cite{bentivogli_cascade} порівняли якість перекладу
у каскадних та end-to-end системах і виявили, що каскадний ефект є
найбільш критичним для сегментів із високою щільністю іменованих
сутностей (власних назв, географічних назв), що узгоджується з
результатами нашого дослідження. Іранзо-Санчес та
ін.~\cite{iranzo_cascade} підтвердили ці спостереження для задачі
автоматичного субтитрування.

\paragraph{ASR для української мови.}
Якість автоматичного розпізнавання української мови значно покращилась
із появою мультимовних моделей. Модель Whisper~\cite{whisper},
натренована на 680\,000 годин слабко-розміченого аудіо, демонструє
WER 10--15\% для українського мовлення на загальних бенчмарках. Для
порівняння, wav2vec~2.0~\cite{wav2vec} потребує fine-tuning на
українських даних і забезпечує порівнянну якість лише за наявності
достатнього обсягу розмічених даних. Розвиток українських мовних
ресурсів та інструментів описано в~\cite{uk_nlp}.

\paragraph{Машинний переклад для мовної пари uk--en.}
Переклад з української на англійську належить до категорії перекладу
з низькоресурсних мов~\cite{low_resource_mt}, що характеризується
обмеженим обсягом паралельних корпусів. Корпус OPUS~\cite{opus_corpus}
є основним джерелом даних для тренування моделей цієї мовної пари.
Спеціалізована модель OPUS-MT~\cite{opus_mt} на архітектурі
MarianNMT~\cite{marian_nmt} досягає BLEU 25--30 на бенчмарку
FLORES~\cite{flores}. Мультимовні моделі NLLB-200~\cite{nllb}
демонструють BLEU 35--38 для пари uk$\to$en. Адаптація MT-моделей
до специфічних доменів (зокрема, політичного дискурсу) залишається
відкритою проблемою~\cite{domain_adaptation}.

\paragraph{LLM як системи машинного перекладу.}
Хенді та ін.~\cite{hendy_llm_mt} провели масштабне оцінювання GPT-моделей
як систем MT та показали, що GPT-4 перевищує спеціалізовані NMT-системи
для більшості мовних пар, зокрема за рахунок кращого розуміння контексту
та здатності до disambiguation. Чжу та ін.~\cite{llm_mt} підтвердили
ці результати для мультимовного перекладу та визначили, що перевага LLM
є найбільш виразною для доменів з високою щільністю іменованих сутностей.

\paragraph{Оцінювання якості перекладу.}
Стандартні автоматичні метрики \\(BLEU~\cite{papineni_bleu},
BERTScore~\cite{bert_score}, METEOR~\cite{banerjee_meteor}) широко
використовуються для оцінювання MT, проте мають відомі
обмеження~\cite{post_bleu}. Для детального аналізу помилок застосовують
фреймворк MQM (Multidimensional Quality Metrics)~\cite{mqm}, що
дозволяє класифікувати помилки за типом та серйозністю. Метрика
WER~\cite{wer_metric} залишається стандартом для оцінювання ASR,
хоча для морфологічно багатих мов додатково використовують CER.

Таким чином, існуючі дослідження демонструють важливість каскадного
ефекту та потенціал LLM для перекладу, проте комплексних досліджень
якості перекладу саме для пари українська--англійська з аналізом
каскадного ефекту та порівнянням із LLM до цього не проводилось.

\subsection{Висновки до розділу 2}

У другому розділі наведено основні терміни та означення: формалізацію
задач ASR та MT, опис каскадної архітектури перекладу мовлення, метрики
оцінювання семантичної точності (BLEU, BERTScore, METEOR), означення
міжфразової когерентності та класифікацію помилок перекладу. Проведено
огляд літератури щодо каскадного ефекту поширення помилок, ASR для
української мови, машинного перекладу для пари uk--en та використання
LLM як MT-систем. Огляд показує, що комплексне дослідження якості
перекладу для пари uk--en з аналізом каскадного ефекту та порівнянням
із LLM до цього не проводилось.

\clearpage
